% !TeX root = ../main.tex
\summary

本文围绕面向数据库 XML 类型的压缩与索引优化方法，研究编码阶段产生的路径、文本和值标识如何继续服务于查询处理和索引构建。在梳理 XML 压缩、压缩域查询以及路径和内容索引研究的基础上，本文提出以中间标识复用为核心的压缩表示方法，给出精确内容、路径存在和路径取值的查询处理过程，并通过 PostgreSQL 原型进行正确性与基本性能验证。

本文表示将共享字典 ID、Varint 内容流、排序路径索引、Bloom 元数据、长值内联区和前序结构流纳入统一数据契约。查询方法在这些结构上分别完成候选生成、否定过滤和精确确认，并使 GIN 操作类从压缩值提取与查询同语义的 token 键。原型当前实现的解析与查询能力界定了实验验证范围，而不是本文方法的研究主题。主要研究结论如下。

第一，结构感知压缩与压缩编码复用是本文方法的核心。本文不把 XML 仅视为一段文本进行通用压缩，而是识别其层次、路径和内容特征，对重复出现的路径与内容进行字典编码，减少原始标签和字符串的重复存储。编码过程中形成的路径 ID 和 token ID 同时服务于压缩表示、查询处理与 GIN 键抽取，使压缩产生的中间数据得到复用。必要的结构信息由结构流保留，以支持解码和结构恢复；索引返回的候选结果则通过精确 recheck 验证。结构差分、查询预言机和回归用例表明，这条方法链在已测范围内能够正确实现。

第二，本文追求的不是单列空间最小，而是在支持高效查询的前提下减少存储开销。与以压缩率为主要目标的 zlib、zstd 相比，结构感知表示需要保留供查询使用的编码和辅助信息，因此单列压缩率并非最优；但在本次 10K 数据实验中，其列值仍比原生 XML 小 20.17\%。与支持相同 exact-token 查询的 $\texttt{text[] + GIN}$ 工程基线相比，本文方法的查询延迟基本相当；将一份共享字典全额计入后，测量持久化成本低 34.38\%。这些结论仅适用于本次数据、实现和统计口径。

在一个简化模型中，设重复结构出现 $N$ 次，不同路径数为 $P$；原始表示中每次重复结构平均占 $s$ 字节，编码后每次平均占 $e$ 字节；字典及路径元数据等固定或随路径数变化的开销记为 $H(P)$。则可近似表示为
\[
\frac{S_{\mathrm{enc}}}{S_{\mathrm{raw}}}
\approx \frac{Ne+H(P)}{Ns}
=\frac{e}{s}+\frac{H(P)}{Ns}.
\]
该关系说明：当数据增长主要来自相似结构重复、$P$ 相对 $N$ 增长较慢，且 $s$、$e$ 与元数据函数 $H(P)$ 的假设大体稳定时，元数据摊销项 $H(P)/(Ns)$ 会随 $N$ 增大而减小，压缩比趋向由 $e/s$ 决定。数据越密集、结构越相似，若它们具体体现为重复结构数量增加而不同路径种类没有同比增加，则同样会摊薄路径字典及辅助结构的单位开销。该式是基于简化假设的趋势分析，不是对任意 XML 数据分布的普遍定理；实际效果仍受文本值分布、长 token 比例、结构流和 Bloom 参数影响。

第三，压缩表示能够支持查询，而不必先还原完整 XML。在本文相同 exact-token 查询语义下，本文方法与 $\texttt{text[] + GIN}$ 工程基线延迟基本持平（27.032 ms 与 27.020 ms）。编码阶段形成的 token ID 可用于 GIN 键抽取和查询匹配；对单个文档，Bloom Filter 能在多数阴性场景快速判定 token 不存在，可能存在的结果仍需经过内容流精确检查。因此，本文的查询价值不在于证明比该基线更快，而是在获得一定空间节省的同时，保持与本次测量中高效索引方案相当的查询性能。
